\documentclass[a4paper]{book}

%%% INICIO DEL PREÁMBULO %%%

\usepackage[utf8]{inputenc}
\usepackage[greek,spanish,es-tabla,es-nodecimaldot,es-noindentfirst]{babel}
\usepackage{babelbib}
\usepackage{nccmath}
\usepackage{amsthm}
\usepackage{lipsum}
\usepackage{tcolorbox}
\usepackage[thicklines]{cancel}
\usepackage{mathtools}
\usepackage{amssymb}
\usepackage{amsmath}
\usepackage{caption}
\usepackage{subcaption}
\usepackage{color}
\usepackage{verbatim}
\usepackage{enumerate}
\usepackage{geometry}
\geometry{a4paper,left=35mm,right=35mm,top=15mm,bottom=15mm}
\usepackage{isotope}
\usepackage{maybemath}
\usepackage{upgreek}
\usepackage{wasysym}
\usepackage[italic]{hepparticles}
\usepackage{subdepth}
\usepackage{siunitx}
\sisetup{
	mode 			= text,
	parse-units 	= false
}
\usepackage{physics}
\usepackage{braket}
\usepackage{tensor}
\usepackage{chemformula}
\usepackage{tikz}
\usepackage{url}
\usepackage{listings}
\usepackage{multirow}
\usepackage{multicol}
\usepackage[colorlinks=true]{hyperref}
\hypersetup{
	citecolor = blue,
	linkcolor = blue,
	urlcolor = blue,
	pdfauthor = {Javier Rodrigo López}
}
\usepackage{eso-pic}

% tikz
\usepackage{tikz} \usetikzlibrary{fit,babel,shapes,arrows,patterns,positioning,calc,decorations.pathmorphing,decorations.markings}
\tikzstyle{block} = [draw, fill=white, rectangle,
minimum height=3em, minimum width=6em]
\tikzstyle{sum} = [draw, fill=white, circle, node distance=1cm]
\tikzstyle{input} = [coordinate]
\tikzstyle{output} = [coordinate]
\tikzstyle{pinstyle} = [pin edge={to-,thin,black}]
\tikzset{
	block/.style = {draw, fill=white, rectangle, minimum height=3em, minimum width=3em},
	tmp/.style  = {coordinate},
	sum/.style= {draw, fill=white, circle, node distance=1cm},
	input/.style = {coordinate},
	output/.style= {coordinate},
	pinstyle/.style = {pin edge={to-,thin,black}}
}

\usepackage[oldvoltagedirection]{circuitikz}
\usepackage{pdflscape}

% Títulos
\usepackage{titlesec}
\titleformat{\section}{\normalfont\Large\bfseries}{\thesection}{1em}{}[{\titlerule[0.8pt]}]
% \renewcommand{\thesubsection}{\arabic{chapter}.\arabic{section}.\Alph{subsection}}
\titleformat{\subsubsection}{\normalfont\normalsize\bfseries}{\thesubsubsection}{1em}{}[{\titlerule[0.05pt]}]
\titlespacing{\section}{0pt}{2\parskip}{\parskip}
\titlespacing{\subsection}{0pt}{\parskip}{0pt}
\titlespacing{\subsubsection}{0pt}{\parskip}{0pt}

% Numeración de secciones
\setcounter{tocdepth}{2}
\setcounter{secnumdepth}{2}

% Figuras y descripciones
\renewcommand{\thefigure}{\arabic{figure}}
\renewcommand{\thesubfigure}{\Alph{subfigure}}
\captionsetup[figure]{labelfont={bf},name={Figura},labelsep=period}
\numberwithin{figure}{chapter}
\numberwithin{equation}{chapter}

% Enumerations
\newcounter{myenumi}
\renewcommand{\themyenumi}{\alph{myenumi})}
\newenvironment{myenumerate}{\setlength{\parindent}{0pt}\setcounter{myenumi}{0}\renewcommand{\item}{\par\refstepcounter{myenumi}\makebox[1.3em][l]{\themyenumi}}}{\par\bigskip\noindent\ignorespacesafterend}

% Own environments
\newenvironment{nota}{\underline{\textbf{NOTA:}} }{}
\newenvironment{caja}{\begin{tcolorbox}[colback = white, sharp corners, boxrule = 1 pt]}{\end{tcolorbox}}
\newtheorem*{conclusion}{Conclusión}
\newtheorem{teorema}{Teorema}
\newtheorem{definicion}{Definición}

% Para una bonita portada
\usepackage{wallpaper}
\usepackage{titling}
\usepackage{fancyhdr}
\pagestyle{fancy}
\setlength{\droptitle}{-10cm}
\renewcommand{\chaptermark}[1]{%
	\markboth{#1}{}}
\renewcommand{\sectionmark}[1]{%
	\markright{}}
\fancyhf{}
\fancyhead[LE,RO]{\bfseries\thepage} \fancyhead[LO]{\bfseries\rightmark} \fancyhead[RE]{\bfseries\leftmark} \renewcommand{\headrulewidth}{0pt} \renewcommand{\footrulewidth}{0pt} \addtolength{\headheight}{15pt}
\fancypagestyle{plain}{%
	\fancyhead{}
	\renewcommand{\headrulewidth}{0pt}
}

% Organización del texto
\newcommand{\formula}[1]{\vspace{13 pt}\noindent \textbf{\underline{#1}}}
\newcommand{\subtext}[1]{_{\text{#1}}}

% Unidades y utilidades varias
\renewcommand{\S}{\operatorname{S}}
\newcommand{\dB}{\operatorname{dB}}
\newcommand{\dBW}{\operatorname{dBW}}
\newcommand{\dBm}{\operatorname{dBm}}
\newcommand{\Hz}{\operatorname{Hz}}
\newcommand{\s}{\operatorname{s}}
\newcommand{\A}{\operatorname{A}}
\newcommand{\V}{\operatorname{V}}
\newcommand{\ohm}{\,\Omega}
\newcommand{\Pa}{\operatorname{Pa}}
\newcommand{\W}{\operatorname{W}}
\newcommand{\I}{\operatorname{I}}
\newcommand{\C}{\operatorname{C}}
\newcommand{\K}{\operatorname{K}}
\newcommand{\m}{\operatorname{m}}
\newcommand{\mm}{\operatorname{mm}}
\newcommand{\rad}{\operatorname{rad}}
\newcommand{\mol}{\operatorname{mol}}
\newcommand{\J}{\operatorname{J}}
\newcommand{\kg}{\operatorname{kg}}
\newcommand{\incremento}{\Delta}
\newcommand{\psus}{\, \ldots \,}
\newcommand{\mcm}{\operatorname{mcm}}
\newcommand{\MCD}{\operatorname{MCD}}
\renewcommand{\sin}{\sen}
\renewcommand{\arcsin}{\arcsen}
\renewcommand{\arctan}{\arctg}
\renewcommand{\min}{\operatorname{mín}}

\DeclarePairedDelimiter\evaluat{.}{\rvert}

% Vectores
\usepackage[c]{esvect}
\renewcommand{\vec}[1]{\vv{{#1}}}
\newcommand{\proy}[2]{\operatorname{proy}_{\vec{#2}}\vec{#1}}
\newcommand{\antiparallel}{\downharpoonleft \! \upharpoonright}
\newcommand{\parallelvec}{\upharpoonleft \! \upharpoonright}

% Espaciado
\usepackage{enumitem}
\setlist{before={\parskip=3pt}, after=\vspace{\baselineskip}}
\setlength{\parindent}{0pt}
\setlength{\parskip}{0.5em}

% Estadística
\DeclareMathOperator{\Var}{Var}
\DeclareMathOperator{\Cov}{Cov}
\renewcommand{\var}{\sigma ^2}
\DeclareMathOperator{\B}{B}
\DeclareMathOperator{\BN}{BN}
\DeclareMathOperator{\Geo}{Geo}
\DeclareMathOperator{\Poisson}{Poisson}
\DeclareMathOperator{\U}{U}
\DeclareMathOperator{\Exp}{Exp}
\DeclareMathOperator{\N}{N}
\DeclareMathOperator{\Mult}{Mult}
\newcommand{\TF}[1]{\mathrm{TF} \left\lbrace \left. #1 \right\rbrace \right.}
\newcommand{\probCond}[2]{P \left( #1 \: \middle\vert\:  #2 \right) }

% Electromagnetismo y Ondas
\newcommand{\errorGrave}{\textbf{FG!!!}}
\newcommand{\mas}{M.A.S.}
\newcommand{\mcu}{M.C.U.}
\newcommand{\ed}{E.D.}
\newcommand{\edmas}{E.D. del M.A.S.}
\usepackage{esint}

% Señales y Sistemas
\renewcommand{\H}{H}

% Circled number
\newcommand{\circledNumber}[1]{\raisebox{.9pt}{\textcircled{\raisebox{-.9pt}{#1}}}}

% Footnotes
% \renewcommand{\thefootnote}{\fnsymbol{footnote}}

% Ejemplo
\newcounter{elejemplo}
\newcommand{\ejemplo}[2]{
	\refstepcounter{elejemplo}
	\begin{center}
		\fbox{\begin{minipage}{0.85\linewidth}
			\textbf{Ejemplo \arabic{elejemplo}.} #1
			\begin{center}
				\underline{\textbf{Solución}}
			\end{center}
			#2
		\end{minipage}}
	\end{center}
}

% Repeticiones
\usepackage{forloop}
\newcommand{\repvec}[3]{
	\foreach \uwu in {1,...,#2}
		{\vec{#1}_{\uwu} ,}
	\, \ldots \, , \vec{#1}_{#3}
}
\newcommand{\rep}[3]{
	\foreach \uwu in {1,...,#2}
		{#1_{\uwu} ,}
	\, \ldots \, , #1_{#3}
}
\newcommand{\repinf}[3]{
	\foreach \uwu in {#2,...,#3}
		{#1_{\uwu} ,}
	\, \ldots
}

%%% FIN DEL PREÁMBULO %%% % Se incluye el preámbulo

% Título y portada
\title{\Huge Electrónica II\\
\Large Apuntes de clase}
\author{Javier Rodrigo López \thanks{E-mail: \href{mailto:javiolonchelo@gmail.com}{\texttt{javiolonchelo@gmail.com}}}} 
\date{\today}

%%% INICIO DEL DOCUMENTO %%%
\begin{document}


\setlength{\wpYoffset}{-5 cm}
\ThisCenterWallPaper{0.7}{./Imágenes/El Bosco.jpg}
\maketitle


% Marca de agua
\AddToShipoutPictureFG{
	\begin{tikzpicture}[overlay,remember picture]
		\path (current page.south west) -- (current page.north east)
		node[midway,scale=8,color=lightgray,sloped,opacity=0.05] {Javier Rodrigo López};
	\end{tikzpicture}
}

% Logotipos UPM y ETSIST
\begin{figure}[t!]
	\centering
	\begin{subfigure}[b]{0.65\linewidth}
		\includegraphics[width=\linewidth]{../../Archivos comunes/upm_logo.png}
	\end{subfigure}
	\begin{subfigure}[b]{0.25\linewidth}
		\includegraphics[width=\linewidth]{../../Archivos comunes/etsist_logo.png}
	\end{subfigure}
\end{figure}


\newpage

\phantomsection
\setlength{\parskip}{0.5em}

\addcontentsline{toc}{section}{Introducción}
\section*{Introducción}
Esta pequeña recopilación de fórmulas, teoremas y demás apuntes de teoría ha sido elaborada durante el primer semestre del curso 2019-2020, en la escuela \href{https://www.etsist.upm.es/}{\textbf{ETSIST}} de la \href{http://www.upm.es/}{\textbf{UPM}} por Javier Rodrigo López, alumno de 2º curso de Ingeniería de Sonido e Imagen.

Muchas gracias al profesor \href{https://www.euitt.upm.es/escuela/directorio?departamento=DTE&idTrabajador=c0874849914b0380a54a53f80d4d60f1}{\textbf{José Antonio Herrera Camacho}} por hacer una labor excelente como docente y despertar en mí el interés por la Electrónica. Una gran parte de estos apuntes existiría si no hubiera asistido a sus clases.
\newpage

\setlength{\parskip}{0em}
\tableofcontents
\setlength{\parskip}{0.5em}


\chapter{Bloque Temático 1}

\section{Codificación de la información}

Para convertir una señal analógica en digital, tenemos que aplicarle una transformación denominada \textbf{codificación}.

La codificación
\section{Codificación de números}
\subsection{Código BN (Binario Natural)}

La fórmula es \[n\geq \log_2{m}\]
Siendo $n$ el número de bits y $m$ el número de objetos. Aunque la mayoría de las veces

El rango de números naturales que se pueden codificar con $n$ bits es:
\[\left( 0, 2^{n}-1 \right)\]

Existen distintas opciones de asignación de códigos. \begin{itemize}
	\item A cada elemento del conjunto debe corresponderle un único código de 0s y 1s.
	\item Se debe utilizar el \textbf{menor número de bits} posible.
	\item La asignación debe favorecer las operaciones de codificar y descodificar.
\end{itemize}
\subsection{Código BCD (Binary Coded Decimal)}
Representa cada una de las cifras decimales de un número expresado en decimal con 4 bits.

Un número decimal de $n$ cifras necesita $4n$ bits para expresarlo en BCD.

El BCD se utiliza mucho para interfaces con personas, ya que cada \textit{nibble} representa
\section{Aritmética Binaria}

\subsection{Suma de números}
De toda la vida. Aquí van algunos ejemplos...

Puede sobrepasar el número de bits con el que se trabaja y se produce un \textbf{acarreo} o \textbf{carry}.

\subsection{Resta de números}
De toda la vida. Aquí van algunos ejemplos...

Puede producirse un error por no utilizar suficientes bits y que el resultado no sea fiel a la operación. A esto se le denomina \textbf{overflow}.

\section{Ejercicios sobre codificación y Aritmética Binaria}


\section{Álgebra de Boole}
\subsection{Teoremas básicos}

\begin{itemize}
	\item \textbf{Teorema de idempotencia}:
	      \begin{multicols}{2}
		      \[\boxed{a + a = a }\]
		      \[\boxed{ a \cdot  a = a}\]
	      \end{multicols}
	\item \textbf{Teorema de absorción}:
	      \begin{multicols}{2}
		      \[\boxed{a + ab = a}\]
		      \[\boxed{ a \cdot  (a+b) = a}\]
	      \end{multicols}
	\item \textbf{Teorema de adyacencia}:
	      \begin{multicols}{2}
		      \[\boxed{ab + a\overline{b} = a}\]
		      \[\boxed{ (a+b)(a+\overline{b}) = a}\]
	      \end{multicols}
	\item \textbf{Teorema de simplificación}:
	      \begin{multicols}{2}
		      \[\boxed{a + \overline{a}b = a + b}\]
		      \[\boxed{ a (\overline{a}+b) = ab}\]
	      \end{multicols}
	\item \textbf{Teorema de De Morgan}:
	      \begin{multicols}{2}
		      \[\boxed{\overline{a_0 \cdot a_1\cdot \ldots \cdot a_n} = \overline{a_0} + \overline{a_1} + \ldots + \overline{a_n}}\]
		      \[\boxed{ \overline{a_0 + a_1+ \ldots + a_n} = \overline{a_0} \cdot  \overline{a_1} \cdot  \ldots \cdot  \overline{a_n} }\]
	      \end{multicols}
\end{itemize}

\section{Cronogramas}
Un cronograma es un diagrama donde se representan las entradas, nodos intermedios y salidas en función del tiempo.

Para esta parte del temario, y a falta de laboratorio en esta asignatura, intentaremos utilizar la aplicación Quartus. Esta herramienta CAD nos permitirá trabajar con puertas lógicas.


\section{Sistemas combinacionales complejos}
La complejidad de un circuito combinacional depende sobre todo del \textbf{número de entradas}, ya que con $N$ entradas la tabla de verdad tendrá $2^N$ filas.

Hemos encontrado casos en los que los circuitos tenían muy poquitos unos o muy poquitos ceros en su tabla de verdad, o que se podían escribir de forma resumida.

\subsection{Divide y vencerás}
La estrategia principal consiste en dividir el circuito en otros más pequeños que, al agruparlos, sean fácilmente descriptibles.

Por ejemplo, un sumador de $N$ bits. Podemos hacer un sumador de 1 bit y con

\chapter{Bloque Temático 2}
\section{Arquitecturas digitales I}
\subsection{Sistemas basados en microprocesador}
\subsection{Sistemas cableados}
Pros: \begin{itemize}
	\item Velocidad y eficiencia.
\end{itemize}
Contras: \begin{itemize}
	\item Complejidad del hardware.
	\item Versatilidad.
\end{itemize}

\section{Conceptos básicos}
\subsection{Nivel logico}
La representación de ceros y unos se define por intervalos de tensiones.

Es decir, si el nivel alto se corresponde con una tensión $V_{CC}$

\subsection{Células lógicas}
Son circuitos electrónicos que realizan operaciones lógicas.

\subsubsection{Nodos lógicos}
Un nodo lógico es una interconexión entre células lógicas.

\subsection{Modelo lógico de las células}
Células estándar: Función lógica prefijada en el diseño de la célula. Por ejemplo, una NOR de dos entradas, construida con transistores CMOS.


\section{Tecnologías}
\subsection{Circuitos cableados}
\subsubsection{SSI. Ciruitos de baja escala de integracíón}
Tecnología de los años 60. Métodos manuales de diseño. Aplicación marginal. Lógica auxiliar (Glue logic)
\subsubsection{Ciruitos de media escala de integracíón}
MSI.
Como mucho, 100 transistores.
\subsubsection{ASIC. Ciruitos de media escala de integracíón}
Tecnología VLSI (Very Large Scale Integration).

\subsubsection{Lógica configurable}
Tecnología VLSI.
Funcionamiento lógico configurable mediante la descarga de un fichero.
\section{Arquitecturas digitales II}

\subsection{Estructura de un sistema digital}

PCB (Printed Circuit Board). Es básicamente la placa donde se colocan los componentes

Alimentación
Interfaces
Disipadores

\section{Tecnología II}
\subsection{Introducción al modelado de CIs}
Necesitamos modelos simplificados para abordar el diseño de sistemas muy complejos.
Hay tres modelos:
\begin{enumerate}
	\item Lógico. Función lógica e interfaz.
	\item Eléctrico. Características eléctricas de la interfaz.
	\item Dinámico.
\end{enumerate}
\subsubsection{Modelo lógico}
Decribe el funcionamiento y la interfaz del chip.
\subsubsection{Modelo eléctrico}
Es muy importante para conectarlo

\chapter{Bloque Temático 3}


\section{Introducción a los circuitos secuenciales}
\section{Cronogramas funcionales de circuitos de flip-flops}
\section{Registros}
\section{Diseño de autómatas}
Un \textbf{autómata de Moore} es un autómata en el que la salida depende únicamente del estado actual de la memoria.

Diagrama de estados para autómatas de Moore:

Un \textbf{autómata de Mealy} es un autómata en el que la salida depende del estado actual de la memoria y de las entradas.

Diagrama de estados para autómatas de Mealy:

\section{Contadores}
\section{Metodología completa de diseño de sistemas}




% BIBLIOGRAFÍA
%\newpage
%\phantomsection
%\label{sec:bibliografia_final}
%\renewcommand{\refname}{Bibliografía} % Para cambiar el título al incluir una librería de BibTeX
%\addcontentsline{toc}{section}{Bibliografía}
%\bibliography{biblio_Tex}
%\bibliographystyle{bababbrv} 

\end{document}
